\begin{table}[htbp]
\caption{Requirements carried from the empirical studies into validation. This synthesis connects the cases; it does not report additional experiments.}
\label{tab:part3-validation}
\centering\small
\renewcommand{\arraystretch}{1.17}
\begin{tabularx}{\linewidth}{@{}>{\raggedright\arraybackslash}p{23mm}>{\raggedright\arraybackslash}X>{\raggedright\arraybackslash}X@{}}
\toprule
\textbf{Case} & \textbf{What the evidence supports} & \textbf{Next validation question}\\
\midrule
\textbf{Recognition}\newline Course assessment & Agreement with specified judgments, documented moderation and correction, and sensitivity of individual positions. & Does the procedure represent work consistently, permit effective correction, and support decisions in another setting?\\[5pt]
\textbf{Continuity}\newline Archived projects & Reasons for archival, observed transition practices, historical trajectories, and a retrospective classification structure. & Do defined functions and responsibilities continue after transition, over the intended horizon?\\[5pt]
\textbf{Governance}\newline Redis and Valkey & Changes in observed participation, selected organizational contributions, and collaboration across related projects. & Are attribution and aggregation defensible, alternatives comparable, and affected parties' outcomes represented?\\
\bottomrule
\end{tabularx}
\tablenote{Evidence is developed in Chapters~\bookref{ch:seven}{7}--\bookref{ch:governance}{9}. Prospective, causal, and normative claims retain different requirements. The validation questions are the book's synthesis of those requirements.}
\end{table}
